% Part I: foundations (sections 1-3)
\addcontentsline{toc}{part}{Part I \quad Foundations}
\section{The toolbox}\label{sec:toolbox}

\subsection{Slopes, curvature, and the second difference}

A function $f$ assigns a number $f(x)$ to every number $x$. Its \emph{derivative} $f'(x)$ is the slope of
the graph at $x$: for a small step $h$,
\[
  f'(x)\approx\frac{f(x+h)-f(x)}{h}.
\]
The second derivative $f''(x)$ is the slope of the slope, which measures curvature: positive where the graph
bends upwards like a cup, negative where it bends down like a cap. Taylor's formula says that near $x$ every
smooth function looks like a parabola,
\begin{equation}\label{eq:taylor}
  f(x+h)=f(x)+h\,f'(x)+\tfrac{1}{2}h^2f''(x)+\tfrac16h^3f'''(x)+\dots
\end{equation}

Write \eqref{eq:taylor} once with $+h$ and once with $-h$ and add. The odd powers cancel, and
\begin{equation}\label{eq:seconddiff}
  \frac{f(x+h)-2f(x)+f(x-h)}{h^2}=f''(x)+O(h^2).
\end{equation}
The left side is the \emph{second difference}. Rewrite it as $\frac{2}{h^2}\big[\tfrac12(f(x+h)+f(x-h))-f(x)\big]$:
it is the gap between the average of the two neighbours and the value at the point. It is positive when the
point sits below its neighbours, negative when it sits above.

\begin{keyidea}
The second difference \eqref{eq:seconddiff} is all a \emph{Laplacian} is in one dimension. When physicists
write $\partial_x^2\phi$ or $\nabla^2\phi$ for a quantity $\phi$ that varies along a line, they mean ``how far
$\phi$ at a point is below the average of its neighbours'', in the limit of close neighbours. In more
dimensions one adds the same quantity for each direction. We will meet it as the force of the springs that
hold a string together (\S\ref{sec:allencahn}).
\end{keyidea}

A function of several variables $f(x_1,\dots,x_n)$ has \emph{partial derivatives} $\partial f/\partial x_i$:
the slope in direction $i$ with the others held fixed. Collected in a vector they form the \emph{gradient}
$\nabla f$, which points in the direction in which $f$ increases fastest.

\subsection{Integrals as sums, and the grid}

The integral $\int_a^b f(x)\dd x$ is the area under the graph. Chop $[a,b]$ into $N_g-1$ intervals of width
$\delta$ with points $u_1,\dots,u_{N_g}$; then
\begin{equation}\label{eq:trapezoid}
  \int_a^b f(x)\dd x\approx\sum_{k=1}^{N_g}w_k\,f(u_k),\qquad w_k=\delta\ \text{except}\ w_1=w_{N_g}=\delta/2 .
\end{equation}
This is the \emph{trapezoid rule}; its error shrinks like $\delta^2$. Every integral in this project is done
this way: on $N_g=401$ points spread over $[-8,8]$, so $\delta=0.04$.

\subsection{Exponentials, logarithms, and why we compute in logs}

The exponential turns sums into products, $e^{a+b}=e^ae^b$, and the logarithm turns products back into sums.
Probabilities of long sequences are products of many numbers smaller than one, which underflow a computer
(in double precision $e^{-800}$ is exactly zero). So we store logarithms, and to add numbers stored as logs
we use the \emph{log-sum-exp} identity
\[
  \log\sum_k e^{y_k}=m+\log\sum_k e^{y_k-m},\qquad m=\max_k y_k,
\]
which never overflows because every exponent on the right is at most zero.

\subsection{Vectors, matrices, eigenvectors}

A vector is a list of numbers, a matrix a table. The product of two matrices is
\begin{equation}\label{eq:matprod}
  (AB)_{kj}=\sum_m A_{km}B_{mj},
\end{equation}
a sum over an intermediate index $m$. Remember this shape: ``sum over the state in the middle'' is the one
operation that belief propagation repeats (\S\ref{sec:bp}). The transpose $A\T$ swaps rows and columns; $A$ is
\emph{symmetric} if $A=A\T$. The inverse $A^{-1}$ undoes $A$.

An \emph{eigenvector} $v$ of a square matrix $A$ is a direction that $A$ only stretches: $Av=\lambda v$, with
\emph{eigenvalue} $\lambda$. Two facts carry most of the weight in this primer.
\begin{itemize}
\item A symmetric matrix has real eigenvalues $\lambda_0\ge\lambda_1\ge\dots$ and eigenvectors $v_0,v_1,\dots$ that
  are mutually perpendicular and of unit length. It can be rebuilt from them, $A=\sum_n\lambda_nv_nv_n\T$, and so
  $A^N=\sum_n\lambda_n^Nv_nv_n\T$. For large $N$ the largest $|\lambda_n|$ dominates: \emph{repeated multiplication
  forgets everything except the leading eigenvector}. That is the power method, and it is the reason a long
  chain ``forgets'' its start.
\item \emph{Perron--Frobenius}: if every entry of $A$ is positive, the largest eigenvalue $\lambda_0$ is positive
  and unrepeated, and its eigenvector can be chosen with every entry positive.
\end{itemize}
A symmetric $A$ is \emph{positive definite} if $v\T Av>0$ for every nonzero $v$, equivalently if all its
eigenvalues are positive. Covariance matrices are of this kind.

\subsection{The Gaussian integral and completing the square}

The single most useful integral is
\begin{equation}\label{eq:gaussint}
  \int_{-\infty}^{\infty}e^{-x^2/2}\dd x=\sqrt{2\pi}.
\end{equation}
(Square it, read the square as an integral over the plane, and switch to polar coordinates; the angular
integral gives $2\pi$ and the radial one gives $1$.) The single most useful trick is \emph{completing the
square}:
\begin{equation}\label{eq:complete}
  -\tfrac{A}{2}x^2+Bx=-\tfrac{A}{2}\Big(x-\tfrac{B}{A}\Big)^2+\tfrac{B^2}{2A}\qquad(A>0).
\end{equation}
It says that any exponent that is quadratic in $x$ describes a bell curve centred at $B/A$ with width
$1/\sqrt A$. Products of Gaussian bumps are Gaussian bumps because their exponents add, and the sum of two
quadratics is a quadratic. The whole of the Kalman filter is this sentence applied repeatedly.

\paragraph{Check yourself.}
\begin{enumerate}[label=\thesection.\arabic*,leftmargin=2.2em]
\item Show that the second difference of $f(x)=x^3$ equals $6x$ exactly, for every $h$. Why is there no error term?
\item The matrix $A=\begin{pmatrix}2&1\\1&2\end{pmatrix}$ has positive entries. Find its eigenvalues and eigenvectors, and say in which
  direction $A^Nv$ points for large $N$, whatever the starting vector $v$ (unless $v$ is special).
\end{enumerate}

\section{Probability from scratch}\label{sec:prob}

\subsection{Random variables, densities, expectations}

A random quantity $a$ is described by its \emph{density} $p(a)$: the probability that $a$ lands in a tiny
interval $[u,u+\dd u]$ is $p(u)\dd u$, and $\int p=1$. The \emph{expectation} of any function of $a$ is the
weighted average
\[
  \E[f(a)]=\int f(u)\,p(u)\dd u .
\]
The mean is $\mu=\E[a]$, the variance $\Var(a)=\E[(a-\mu)^2]$, the standard deviation its square root. For two
quantities the covariance $\Cov(a,b)=\E[(a-\E a)(b-\E b)]$ measures whether they move together, and the
correlation divides it by both standard deviations, landing in $[-1,1]$.

Two random quantities have a \emph{joint} density $p(a,b)$. Integrating one away gives the \emph{marginal},
$p(a)=\int p(a,b)\dd b$. Fixing one gives the \emph{conditional},
\[
  p(a\mid b)=\frac{p(a,b)}{p(b)},
\]
the law of $a$ once $b$ is known. They are independent when $p(a,b)=p(a)p(b)$. Writing the joint both ways,
$p(a\mid x)p(x)=p(x\mid a)p(a)$, gives \emph{Bayes' rule}:
\begin{equation}\label{eq:bayes}
  \underbrace{p(a\mid x)}_{\text{posterior}}=\frac{\overbrace{p(x\mid a)}^{\text{likelihood}}\ \overbrace{p(a)}^{\text{prior}}}{\underbrace{p(x)}_{\text{evidence}}}
  \;\propto\;p(x\mid a)\,p(a).
\end{equation}
The symbol $\propto$ (``proportional to'') means ``equal up to a factor that does not depend on $a$''; the factor
is recovered at the end by requiring the density to integrate to one.

\subsection{The Gaussian, in one and many dimensions}

The Gaussian (normal) density with mean $m$ and variance $s^2$ is
$\N(u;m,s^2)=\exp\!\big(-(u-m)^2/2s^2\big)/\sqrt{2\pi s^2}$; by \eqref{eq:gaussint} it integrates to one. A
weighted sum of independent Gaussians is Gaussian, with means and variances combining as
$\E[\alpha a+\beta z]=\alpha\E a+\beta\E z$ and $\Var(\alpha a+\beta z)=\alpha^2\Var a+\beta^2\Var z$.

A vector $a=(a_1,\dots,a_n)$ is Gaussian with mean vector $m$ and covariance matrix $\Sigma$
($\Sigma_{ij}=\Cov(a_i,a_j)$) if
\begin{equation}\label{eq:mvn}
  p(a)=\frac{1}{\sqrt{(2\pi)^n\det\Sigma}}\exp\!\Big(-\tfrac12(a-m)\T\Sigma^{-1}(a-m)\Big).
\end{equation}
The inverse $\Sigma^{-1}$ is the \emph{precision matrix}, and its zeros have a meaning that will matter: if
$(\Sigma^{-1})_{ij}=0$, then $a_i$ and $a_j$ are independent \emph{once all the other coordinates are known}.

\begin{example}[the Gaussian AR(1) chain has a tridiagonal precision]
Let $a_1\sim\N(0,1)$ and $a_{i+1}=\rho a_i+\varepsilon_i$ with independent $\varepsilon_i\sim\N(0,1-\rho^2)$. The joint
density is $p(a_1)\prod_i\N(a_{i+1};\rho a_i,1-\rho^2)$, whose exponent is
$-\tfrac12a_1^2-\sum_i(a_{i+1}-\rho a_i)^2/2(1-\rho^2)$. Each term couples only neighbours, so the precision
matrix has nonzeros only on the diagonal and the two adjacent diagonals, while the covariance
$\Sigma_{ij}=\rho^{|i-j|}$ is full. \emph{Local in the precision, global in the covariance}: this is the whole
reason inference on chains is cheap.
\end{example}

\subsection{Conditioning a Gaussian on a noisy observation}\label{sec:gausscond}

Suppose $a\sim\N(m,P)$ and we observe $x=\alpha a+\sqrt\Delta\,z$ with independent noise $z\sim\N(0,1)$. By Bayes'
rule the posterior is proportional to the prior times the likelihood:
\[
  p(a\mid x)\propto\exp\!\Big(-\frac{(a-m)^2}{2P}-\frac{(x-\alpha a)^2}{2\Delta}\Big).
\]
The exponent is quadratic in $a$. Collect the coefficient of $a^2$ (that is $-\tfrac12A$ with $A=1/P+\alpha^2/\Delta$)
and of $a$ (that is $B=m/P+\alpha x/\Delta$) and complete the square \eqref{eq:complete}:
\begin{equation}\label{eq:precadd}
  \boxed{\;\frac{1}{P'}=\frac1P+\frac{\alpha^2}{\Delta},\qquad m'=P'\Big(\frac{m}{P}+\frac{\alpha x}{\Delta}\Big)\;}
\end{equation}
The posterior is Gaussian. \emph{Precisions add} (prior precision plus the information in the measurement),
and the posterior mean is a precision-weighted average of what we believed and what we saw. The same result
in ``gain'' form, which is how the Kalman filter is usually written, is
\begin{equation}\label{eq:gain}
  m'=m+k\,(x-\alpha m),\qquad P'=(1-k\alpha)P,\qquad k=\frac{\alpha P}{\alpha^2P+\Delta}.
\end{equation}
The \emph{innovation} $x-\alpha m$ is the surprise in the measurement; the gain $k$ says how much of it to
believe. (Check: $(1-k\alpha)P=P\Delta/(\alpha^2P+\Delta)$, which is the reciprocal of $1/P+\alpha^2/\Delta$.)

For vectors the statement is the same with matrices. If $(a,x)$ is jointly Gaussian with blocks
$\Sigma_{aa},\Sigma_{ax},\Sigma_{xx}$, then
\begin{equation}\label{eq:condmvn}
  \E[a\mid x]=m_a+\Sigma_{ax}\Sigma_{xx}^{-1}(x-m_x),\qquad \Cov(a\mid x)=\Sigma_{aa}-\Sigma_{ax}\Sigma_{xx}^{-1}\Sigma_{xa}.
\end{equation}
Notice that the conditional mean is a \emph{linear} function of $x$. This is special to Gaussians.

\subsection{The best predictor and the best linear predictor}\label{sec:mmse}

Suppose we must guess $a$ from $x$ with a function $g(x)$ and pay the squared error. The best $g$ is the
\emph{conditional mean} $g^\star(x)=\E[a\mid x]$. Proof: for any $g$,
\[
  \E\big[(a-g)^2\big]=\E\big[(a-g^\star)^2\big]+\E\big[(g^\star-g)^2\big],
\]
because the cross term $\E[(a-g^\star)(g^\star-g)]$ vanishes (given $x$, the second factor is fixed and the first
has mean zero). The second term is never negative, so nothing beats $g^\star$. This is the
\emph{minimum mean-square error} (MMSE) estimator, and it is what the word ``denoiser'' means in this primer.

If we restrict ourselves to \emph{linear} guesses $g(x)=Gx$, the best one is
$G=\Cov(a,x)\Var(x)^{-1}$, the LMMSE (linear MMSE) estimator. It depends only on covariances. For a Gaussian
pair it coincides with $\E[a\mid x]$ by \eqref{eq:condmvn}; for anything else, the true conditional mean is
nonlinear and strictly better.

\begin{keyidea}
Every law has a \emph{covariance twin}: the Gaussian with the same means and covariances. The twin's
denoiser is the LMMSE estimator of the true law. The gap between the twin's denoiser and the true one is
exactly what the non-Gaussian structure of the data is worth to a denoiser. Most of this project measures
that gap.
\end{keyidea}

\subsection{Moments, cumulants, kurtosis}\label{sec:cumulants}

For a standardised quantity (mean $0$, variance $1$) the third moment $\E a^3$ is the \emph{skewness}
(asymmetry) and $\E a^4$ the \emph{kurtosis}, which is $3$ for a Gaussian. The \emph{excess kurtosis}
$\E a^4-3$ is positive for heavy tails (a Laplace law has $+3$) and negative for laws with ``shoulders'' or
two humps (the $\phi^4$ marginal at $\Th=0.3$ has $-1.35$).

These are special cases of \emph{cumulants} $\kappa_n$, defined as the coefficients of
$\log\E e^{sa}=\sum_n\kappa_ns^n/n!$. The first is the mean, the second the variance, and for a standardised
law $\kappa_3$ and $\kappa_4$ are skewness and excess kurtosis. Two properties make cumulants the right
language. For independent $a$ and $b$ they add, $\kappa_n(a+b)=\kappa_n(a)+\kappa_n(b)$, and they scale as
$\kappa_n(ca)=c^n\kappa_n(a)$. A Gaussian has $\kappa_n=0$ for every $n\ge3$; this is what ``non-Gaussian'' will
mean.

A first look at the decay that organises the whole project: if $x=\alpha a+\sqrt\Delta z$ with Gaussian $z$
independent of $a$, then for $n\ge3$
\begin{equation}\label{eq:cumdecay}
  \kappa_n(x)=\alpha^n\kappa_n(a)+\Delta^{n/2}\kappa_n(z)=\alpha^n\kappa_n(a).
\end{equation}
With $\alpha=e^{-t}$ (\S\ref{sec:noise}), the $n$-th cumulant fades like $e^{-nt}$: skewness like $e^{-3t}$,
kurtosis like $e^{-4t}$, while the variance is held fixed. Chapter~\ref{sec:hermite} turns this into a
statement about whole functions.

\paragraph{Check yourself.}
\begin{enumerate}[label=\thesection.\arabic*,leftmargin=2.2em]
\item The prior is $a\sim\N(0,1)$ and we observe $x=a+z$ with $z\sim\N(0,1)$ ($\alpha=\Delta=1$). What is the posterior?
\item A standardised law has excess kurtosis $-1.35$. What is the excess kurtosis of $x=e^{-t}a+\sqrt{1-e^{-2t}}\,z$ at $t=1$?
\end{enumerate}

\section{Markov chains}\label{sec:markov}

\subsection{Definition}

A sequence $a_1,a_2,\dots,a_L$ is a \emph{Markov chain} if the next value depends on the past only through the
present: $p(a_{i+1}\mid a_1,\dots,a_i)=p(a_{i+1}\mid a_i)$. The joint density then factorises,
\begin{equation}\label{eq:chain}
  p(a_1,\dots,a_L)=\mu(a_1)\prod_{i=1}^{L-1}K(a_{i+1}\mid a_i),
\end{equation}
into an initial law $\mu$ and a \emph{transition kernel} $K(a'\mid a)$, which for each $a$ is a density in $a'$.
On the grid of \eqref{eq:trapezoid} the kernel becomes a matrix $K_{kj}=K(u_k\mid u_j)$ whose columns satisfy
$\sum_kw_kK_{kj}=1$, and pushing a distribution one step forward is a matrix--vector product:
\[
  p_{i+1}(a')=\int K(a'\mid a)\,p_i(a)\dd a\quad\leadsto\quad p_{i+1}(u_k)=\sum_jK_{kj}\,w_j\,p_i(u_j).
\]

\subsection{The AR(1) chain}

The simplest continuous chain is the first-order autoregression
\begin{equation}\label{eq:ar1}
  a_{i+1}=\rho\,a_i+\varepsilon_i,\qquad \E\varepsilon_i=0,\ \Var\varepsilon_i=q,
\end{equation}
with innovations $\varepsilon_i$ independent of the past. If the variance is to stay at $v$ from site to site,
$v=\rho^2v+q$, so $q=1-\rho^2$ keeps it at $1$. Iterating, $a_{i+k}=\rho^ka_i+(\text{noise independent of }a_i)$,
so the correlation between sites $k$ apart is $\rho^k=e^{-k/\xi}$ with \emph{correlation length}
$\xi=1/\log(1/\rho)$ sites. At the project's frozen $\rho=0.85$, $\xi\approx6.2$ sites.

The innovation need not be Gaussian. A Laplace innovation, or a mixture of Gaussians, gives a chain with
\emph{exactly the same covariance} as the Gaussian one and different higher-order structure. That was the
project's first laboratory: hold the second-order structure fixed and vary the rest.

\subsection{Stationarity and reversibility}

A density $\pi$ is \emph{stationary} if one step of the chain leaves it unchanged,
$\pi(a')=\int K(a'\mid a)\pi(a)\dd a$. On the grid, $\pi$ is an eigenvector of the matrix $KW$ ($W$ the diagonal
of weights) with eigenvalue $1$. For the chains we use it is unique and the chain forgets its start
geometrically fast (the power method again). A chain started from $\pi$ is \emph{stationary}: every site has
law $\pi$.

The chain is \emph{reversible} if it satisfies \emph{detailed balance},
\begin{equation}\label{eq:detbal}
  \pi(a)K(a'\mid a)=\pi(a')K(a\mid a'):
\end{equation}
the flow of probability from $a$ to $a'$ equals the flow back. A film of a reversible stationary chain looks
statistically the same played backwards. Detailed balance implies stationarity (integrate both sides over $a$).

\begin{inproject}
The index $i$ need not be time. A snapshot of a one-dimensional field (the value of a physical quantity at
each of $L$ equally spaced points along a line) is a sequence $\phi_1,\dots,\phi_L$. If the physics couples only
neighbouring points, this sequence is a Markov chain \emph{in space}: the value at site $i+1$ depends on the
values to its left only through site $i$. Chapters~\ref{sec:allencahn} and~\ref{sec:transfer} show that the
stochastic Allen--Cahn field is exactly of this kind, and build its kernel $K$.
\end{inproject}

\paragraph{Check yourself.}
\begin{enumerate}[label=\thesection.\arabic*,leftmargin=2.2em]
\item For the AR(1) chain at $\rho=0.85$, what is the correlation between sites ten apart?
\item Show that detailed balance \eqref{eq:detbal} implies that $\pi$ is stationary.
\end{enumerate}
